\subsection{Real descent and the range of the realization}

\begin{result}{C.1.1}{Real Hodge realization on abelian motives}{Proved}
Let $\mathcal C_{\R}^{\mathrm{ab}}$ be the rational tensor category of motives for absolute Hodge cycles over $\R$ generated by abelian varieties and Artin motives. Its realization
\[
 M\longmapsto (H_{\B}(M),f_{\infty,M})
\]
is fully faithful into polarizable rational Hodge structures with a rational involution satisfying $f(H^{p,q})=H^{q,p}$. The essential image of the subcategory consisting of $h^1$ of abelian varieties, up to isogeny, consists exactly of such pairs of Hodge type $\{(1,0),(0,1)\}$. The essential image of $\mathcal C_{\R}^{\mathrm{ab}}$ is the rigid, additive, idempotent-complete tensor subcategory generated by these pairs and the sign Artin line.
\end{result}
\begin{proof}
After base change to $\C$, a morphism between the indicated tensor constructions is a Hodge class on a product of abelian varieties, with a Tate twist and projector. Deligne's theorem makes every such class absolute; we use the form in \cite[Theorem 58]{CharlesSchnell}. Commuting with $f_\infty$ is precisely invariance under the descent action of $\Gal(\C/\R)$. An invariant absolute Hodge class is a morphism over $\R$ by the definition of this category. Faithfulness follows from the faithful Betti realization.

For the degree-one image, we use Riemann's theorem in \cite[Theorem 6.20 and Proposition 6.21]{DM1982}. Here is the additional real-descent step. Dualize to homology and choose a rational lattice $\Lambda_0$ in the underlying rational vector space. The lattice $\Lambda=\Lambda_0+f\Lambda_0$ is $f$-stable. If $E_0$ is a rational Riemann form and $J$ the complex structure, put
\[
 E(x,y)=E_0(x,y)-E_0(fx,fy).
\]
Since $fJ=-Jf$, both terms contribute positively to $E(x,Jx)$, and $f^*E=-E$. Multiplying $E$ by a positive integer makes it integral on $\Lambda$. The complex torus $(\Lambda\otimes\R,J)/\Lambda$ is therefore a polarized abelian variety. The map induced by $f$ is an antiholomorphic group involution. As a holomorphic map to the conjugate variety it is algebraic and satisfies the quadratic descent identity; projective descent gives an abelian variety over $\R$. Its degree-one Hodge pair is the original pair. Taking tensors, duals, sums and projectors now gives the stated generated essential image, since the realization is fully faithful.
\end{proof}

\begin{result}{C.1.2}{The homological-motive version and its scope}{Proved under stated hypotheses}
Assume the Hodge conjecture for all products and Tate twists used to form the homological motives under consideration over $\R$. Their realization into Hodge pairs is fully faithful. This proves the fullness assertion of 18:34 in that scope; it supplies no characterization of the geometric essential image among all polarizable Hodge pairs.
\end{result}
\begin{proof}
A Hodge morphism is represented over $\C$ by an algebraic correspondence $z$ under the stated conjecture. For an $f_\infty$-invariant morphism, the transferred correspondence $\tfrac12\operatorname{Tr}_{\C/\R}(z)$ has the required Betti class over $\R$. The projectors commute with this operation. Equality of morphisms in the homological category is equality of their cohomology classes, which proves faithfulness. The argument only constructs morphisms between already geometric objects, so it makes no claim that an arbitrary Hodge pair is geometric.
\end{proof}

\begin{result}{C.1.3}{Generation by the identity component and real Frobenius}{Proved}
In $\mathcal C_{\R}^{\mathrm{ab}}$, and in the homological category of C.1.2 under its hypothesis, the group attached to any tensor generator is generated by its identity component and $f_\infty$. Its component group is trivial or of order two. A nontrivial component is a translate of the identity component as a $\Q$-scheme.
\end{result}
\begin{proof}
Put $V=H_{\B}(M)$, $f=f_{\infty,M}$ and
$G=\Aut^\otimes(H_{\B}|_{\langle M\rangle})\subset\GL(V)$.
Let $h:\mathbb S_{\R}\to\GL(V_{\R})$ be the Hodge homomorphism,
and let $H\subset\GL(V)$ be the smallest algebraic $\Q$-subgroup
whose scalar extension contains its image. This is the
Mumford--Tate group of $V$. It is connected: the connected group
$h(\mathbb S_{\R})$ contains the identity and is contained in
$H^0_{\R}$, while $H^0$ is defined over $\Q$; minimality gives
$H=H^0$.

The Hodge-pair identity gives
$fh(z)f^{-1}=h(\bar z)$. Since conjugation permutes
$h(\mathbb S_{\R})$, the rational element $f$ normalizes $H$.
As $f^2=1$, the finite union
\[
                         K=H\cup fH
\]
is a closed algebraic $\Q$-subgroup, with $K^0=H$.
The maps $h$ and $f$ act naturally on every tensor construction
and commute with every Hodge-pair morphism. Full faithfulness
in C.1.1 therefore gives $K\subset G$. In every mixed tensor
construction $T(V)$, that same full faithfulness gives
\[
 T(V)^G=\operatorname{Hom}(\mathbf1,T(V))
       =\{t\in T(V):t\text{ has type }(0,0),\ ft=t\}
       =T(V)^K.                                      \tag{C.1.3a}
\]
Here $H$-invariance is precisely invariance under $h$, by its
definition as the smallest algebraic $\Q$-group containing $h$;
for a rational vector, invariance under $h$ is the stated
Hodge-type condition. The assertion includes duals and direct
sums, so it determines Hom spaces between all tensor words.

For completeness, these invariants determine the group in this
setting without a normality assumption. Polarizable rational
Hodge structures are semisimple. If $W\subset X$ is a subobject
of Hodge pairs and $p:X\to W$ is a Hodge retraction, then
$(p+f_Wpf_X)/2$ is a Hodge retraction commuting with $f$.
It is a retraction because both summands restrict to $1_W$;
it commutes with $f$ by $f_X^2=f_W^2=1$. Thus the category of
polarizable Hodge pairs is semisimple, as is the tensor
subcategory in C.1.1. The corresponding $G$ is reductive.
The same semisimplicity of polarizable Hodge structures shows
that $H$ is reductive; its finite extension $K$ is linearly
reductive in characteristic zero, by averaging over $K/H$.
Every finite-dimensional representation of the closed subgroup
$K\subset\GL(V)$ is a subquotient of finite sums of mixed
tensor words in $V$, hence, by this reductivity, a direct
summand. Equation (C.1.3a) says that every $K$-equivariant map
between these words, including every idempotent, is
$G$-equivariant. Restriction $\operatorname{Rep}(G)\to
\operatorname{Rep}(K)$ is therefore fully faithful and
essentially surjective. Tensor reconstruction
\cite[Theorem 2.11]{DM1982} identifies the inclusion $K\hookrightarrow G$
with an isomorphism.

Consequently $G^0=H$ and $G=G^0\cup fG^0$. Its component
group is trivial or the constant group of order two. In the
second case left multiplication by $f\in G(\Q)$ is the
claimed isomorphism of $\Q$-schemes $G^0\to fG^0$.
For the homological category of C.1.2 its stated fullness
supplies the same equality (C.1.3a), and the same proof applies.
\end{proof}

\begin{result}{C.1.4}{An explicit polarized pair outside the abelian image}{Proved}
Put $a=\sqrt[3]{2}>0$, $V=\Q e_1\oplus\Q e_2$, and
\[
 V^{3,0}=\C(e_1+ia e_2),\qquad
 V^{0,3}=\C(e_1-ia e_2),\qquad f=\operatorname{diag}(1,-1).
\]
This is a polarizable rational Hodge structure with a rational anti-Hodge involution. The pair $(V,f)$ is outside the essential image of the abelian-motive category in C.1.1.
\end{result}
\begin{proof}
For any positive real $a$, the displayed two lines are complementary and conjugate, and $f$ exchanges them. In the convention $h(z)|_{V^{p,q}}=z^{-p}\bar z^{-q}$ the Weil operator and the rational alternating coefficient of a polarization are
\[
 C_a=\begin{pmatrix}0&a^{-1}\\-a&0\end{pmatrix},\qquad
 E=\begin{pmatrix}0&-1\\1&0\end{pmatrix},\qquad
 q(x,y)=(2\pi i)^{-3}x^tEy.
\]
The pairing has target $\Q(-3)$, is zero on each of the two equal-type lines, and is therefore a Hodge morphism. Moreover
\[
 EC_a=\operatorname{diag}(a,a^{-1})>0,\qquad
 f^tEf=-E,\qquad fC_af=C_a^{-1}.
\]
Thus $q$ is a polarization, also compatible with the involution on its Tate target. These identities are exact, and apply to both examples in this subsection.

For $a=\sqrt[3]{2}$, put $\tau=ia$. An endomorphism $\begin{psmallmatrix}r&s\\t&u\end{psmallmatrix}$ with rational entries preserves $V^{3,0}$ precisely when
\[
 t+(u-r)\tau-s\tau^2=0.
\]
Here $\tau^2=-\sqrt[3]{4}$ has degree three over $\Q$, and $\tau$ is not real; hence $\tau$ has degree six. The displayed relation forces $s=t=0$ and $u=r$, proving $\End_{\mathrm{HS}}(V)=\Q$. A commutative Mumford--Tate group would have its rational Lie algebra inside this endomorphism algebra. Its dimension is at least two, since its real points contain the nontrivial weight scaling and Hodge circle. Their tangent directions are linearly independent. This contradicts the one-dimensional endomorphism algebra. Thus $V$ is not of CM type.

Totaro's Corollary 4.2, proved from Theorem 4.1 in \cite[p.~18]{TotaroHodge2015}, excludes this $V$ from the tensor category generated by $H^1$ of curves. The realization functor from C.1.1, followed by forgetting $f$, has image in exactly that tensor category: an Artin summand becomes a sum of type $(0,0)$ lines. This explicit composite functor proves the claimed exclusion for the pair.
\end{proof}

\begin{result}{C.1.5}{A fully specified polarized pair outside the geometric image}{Proved}
Fix the following effective enumeration $(T_n)_{n\geq1}$ of deterministic Turing machines. The tape alphabet is $\{0,1,\#,\square\}$ in that order. For each $r\geq1$ use states $0,\ldots,r$, with initial state $0$ and halt state $r$; list the tables for the $4r$ nonhalting state/symbol pairs in lexicographic order, each instruction being an output symbol, a next state, and a move from the ordered list left, stay, right. Order first by $r$. Input $n$ is the binary expansion of $n$ between two $\#$ symbols on an otherwise blank two-sided tape, with the head at the first binary digit. This specifies a single enumeration of all machines of this model. Let
\[
 H=\{n:T_n\text{ halts on input }n\},\qquad
 a_H=1+\sum_{n=1}^{\infty}\frac{\mathbf1_H(n)}{4^n}.
\]
Replace $a$ in C.1.4 by $a_H$. The resulting polarized rational Hodge pair is not the realization of a pure motive defined by algebraic correspondences over $\C$, and consequently is not the realization of such a motive over $\R$. This is an exclusion from the geometric essential image; a description of every pair in that image remains an obligation.
\end{result}
\begin{proof}
We first check the specified real number. It lies in $[1,4/3]$. The set $H$ is not decidable: a machine which, on input $n$, halts exactly when a proposed decider says $T_n(n)$ does not halt gives a contradiction at its own index. If $a_H$ were computable, its base-four digits in $\{0,1\}$ would decide $H$. More explicitly, after the first $n-1$ digits have been recovered, multiply the remaining tail by $4^n$. It lies in $[0,1/3]$ when the $n$th digit is zero, and in $[1,4/3]$ when that digit is one. A computable approximation of error less than $1/6$ separates the intervals. Induction would decide every digit, a contradiction. Thus $a_H$ is noncomputable. The polarization and involution identities of C.1.4 apply because $a_H>0$; no numerical approximation is used.

We spell out why a geometric realization would force this number to be computable. A finite presentation of a pure algebraic-correspondence motive over $\C$ uses finitely many smooth projective varieties, twists, rational projectors and cycles. All of this data descends to a finitely generated extension of $\Qb$. Spreading the defining equations and the finitely many cycles gives, after restricting a nonempty open set, smooth projective families over a smooth integral finite-type $\Qb$-scheme $S$. Cycle actions are horizontal endomorphisms of their rational cohomology local systems. Idempotence, satisfied at the given complex point, holds on its connected component, so their images form a rational variation of Hodge structures. Its Hodge numbers are constant: the image of an idempotent on each Hodge bundle has locally constant rank.

For the asserted rank-two weight-three image one has $F^3=F^2=F^1$. Griffiths transversality gives $\nabla F^3\subset\Omega_S^1\otimes F^2=\Omega_S^1\otimes F^3$. Hence $F^3$ is horizontal; its conjugate is horizontal too. Parallel transport is therefore an isomorphism of rational Hodge structures along any path in the connected component. That component contains a $\Qb$-point: such points are analytically dense in a smooth variety over $\Qb$, as follows in an etale coordinate chart from the analytic inverse-function theorem and density of $\Qb$ in $\C$. A path to that point produces a rational Hodge isomorphism with a motive over $\Qb$. This is the explicit spreading argument behind the gap-in-Hodge-numbers phenomenon noted in \cite[Introduction, pp.~1--2]{TotaroHodge2015}.

For a motive over $\Qb$, choose algebraic de Rham bases, and a rational Betti basis. Comparison entries are periods over algebraic numbers, with integer powers of $2\pi i$ for the twists. Such periods have computable real and imaginary parts by \cite[Theorem 18, p.~11]{YoshinagaPeriods2008}. Algebraic numbers are computable, and the computable complex numbers form a field: reciprocal approximation follows by first finding a positive rational lower bound for the modulus of the nonzero denominator. Consequently the coordinates of the one-dimensional $F^3$ line, and its ratio in every rational basis, are computable. Rational change of basis acts on the ratio by an explicitly invertible fractional-linear transformation with rational coefficients and preserves this property. Applied to the line $\C(1,ia_H)$, this makes $a_H$ computable, contradicting the first paragraph.

This proof uses algebraic projectors in the geometric category. It does not replace them by arbitrary Hodge projectors over an arbitrary complex coefficient field. The forgetful map from real motives to their complex base changes transfers the exclusion in the required direction.
\end{proof}



\subsection{The twisted Betti functor}

\begin{result}{C.2.1}{Tensor structure and the exact torsor}{Proved}
Let $\mathcal C$ be any neutral rational Tannakian category with fibre functor $V$ and a tensor involution $f\in G(\Q)$, $G=\Aut^\otimes(V)$. Write $V(X)=V(X)^+\oplus V(X)^-$ and define
\[
 W(X)=V(X)^+\oplus iV(X)^-\subset V(X)\otimes_\Q\Q(i).
\]
Multiplication in $\Q(i)$ gives a tensor structure on $W$, making it a rational fibre functor. Its comparison torsor is the twist of $G$ by
\[
 a_\sigma=f^{\epsilon(\sigma)},\qquad
 \sigma(i)=(-1)^{\epsilon(\sigma)}i,
 \quad \sigma\in\Gal(\Qb/\Q).
\]
For a field $F$ of characteristic zero containing $\Q$, its $F$-points are, using the inverse matrix of the ambient map $V_{F(i)}\to W_{F(i)}\hookrightarrow V_{F(i)}$,
\[
 \Isom^\otimes(V,W)(F)=
 \{b\in G(F(i)):b^{-1}\bar b=f\}\quad(i\notin F).
\]
For $i\in F$, the torsor is trivial. This specifies its class in $H^1(\Q,G)$ and exactly tests triviality over $\Q$, $\R$ and $\C$.
\end{result}
\begin{proof}
Equivalently $W(X)$ is the fixed space of $f\otimes c$ on $V(X)\otimes\Q(i)$. Descent of finite-dimensional vector spaces proves exactness and faithfulness. The four tensor components multiply with factors $1,i,i,-1$, respectively; hence their images are precisely the four summands of $W(X\otimes Y)$. Associativity, the unit and the symmetry constraints are inherited from the ambient tensor category. Over $\Q(i)$ the inclusions identify $W$ with $V$. Conjugating this identification changes it by $f$, proving the displayed cocycle. The torsor classification used here is \cite[Theorem 3.2]{DM1982}. A trivializing point has an ambient matrix $g\in G(F(i))$ satisfying $f\bar g=g$. With the explicitly chosen inverse coordinate $b=g^{-1}$, this becomes $b^{-1}\bar b=f$. This also specifies the orientation of the displayed nonabelian coboundary.

The additive map $v_++v_-\mapsto v_++iv_-$ has tensor defect $-1$ on $V^-\otimes V^-$: its tensor square gives $-v\otimes w$, whereas its value on $v\otimes w$ is $v\otimes w$. This verifies RH06 without using that particular additive map as a tensor trivialization.
\end{proof}

\begin{result}{C.2.2}{Artin, Tate and elliptic tests}{Proved}
For the sign Artin motive $A$ of $\C/\R$, the twisted Betti functor has no tensor isomorphism to Betti over $\Q$ or $\R$, and has one over $\C$. For the Tate category, whose group is $\Gm$, it is tensor-isomorphic over $\Q$. For the category generated by $h^1(E)$ of a real elliptic curve, it is tensor-isomorphic over $\Q$ in the non-CM case and is not tensor-isomorphic over $\R$ in the CM case. It is tensor-isomorphic over $\C$ in both cases.
\end{result}
\begin{proof}
Choose an Artin basis $a$ with $a^{\otimes2}\mapsto1$. A rational or real tensor isomorphism would send it to $qi a$ and force $-q^2=1$, impossible over either field. The Artin torsor is $\Spec\Q(i)$ with its $\mu_2$-action.

For $T=\Q(-1)$, choose bases $e_n$ on $T^{\otimes n}$ with $e_m e_n=e_{m+n}$. The maps $e_n\mapsto i^ne_n$ give a rational tensor isomorphism to $W$: the codomain is the rational line $i^n\Q e_n$, and there is no finite-order relation on $T$. The cocycle $-1$ in $\Gm$ is the coboundary $i^{-1}\bar i$. Thus the square-class test in Lead H7 applies to the Artin quotient $\mu_2$, not to $\Gm$; these are related by the inclusion $\mu_2\hookrightarrow\Gm$, which sends the nontrivial class represented by $\Q(i)$ to this explicit trivial class.

For non-CM $E$, the geometric group and the real group on $h^1$ are $\GL_2$. Real Frobenius has eigenvalues $1,-1$: it reverses the orientation of the two-dimensional real torus and has determinant $-1$. In a rational eigenbasis, $b=\operatorname{diag}(1,i)$ belongs to $\GL_2(\Q(i))$ and satisfies $b^{-1}\bar b=f$. Its actions on all tensor constructions give the required tensor isomorphism, even though the additive formula in C.2.1 does not do so on all objects.

For CM $E/\R$, the imaginary quadratic field $K=\End^0(E_\C)$ occurs as Hodge endomorphisms of $h^1$. Its trace-zero line inside $\End(h^1)$ has type $(0,0)$ and $f$ acts by $-1$, since conjugation acts as the nontrivial automorphism of $K$. That line is the sign Artin motive $A$, by C.1.1 or by Artin descent. A tensor isomorphism on $\langle h^1(E)\rangle$ would restrict to the impossible one on $A$. Base change to $\C$ trivializes every twist in question.
\end{proof}

\subsection{The exact kernel of real Frobenius}

\begin{result}{C.3.1}{The even-Tate assertion and its replacement}{Refuted; replacement proved}
In $\mathcal C_{\R}^{\mathrm{ab}}$, put $L=A\otimes\Q(-1)$. The full subcategory on which $f_\infty$ acts trivially consists exactly of finite direct sums of $L^{\otimes p}$, $p\in\Z$. The same conclusion holds for homological motives under the Hodge hypothesis of C.1.2. In particular, the assertion on 18:38 that these are sums of tensor powers of $\Q(2)$ is false.
\end{result}
\begin{proof}
The equation $fH^{p,q}=H^{q,p}$ with $f=1$ forces $H^{p,q}=0$ for $p\ne q$, since the Hodge decomposition is direct. Over $\C$, every remaining pure summand of type $(p,p)$ is a sum of $\Q(-p)$: the corresponding Hodge maps are absolute in the abelian category, and algebraic under the named Hodge hypothesis. Descent to $\R$ is an involution on the multiplicity space. Because $\Q(-p)$ contributes $(-1)^p$ to real Frobenius, trivial total Frobenius forces the multiplicity representation to be $A^{\otimes p}$. Conversely $f$ acts on $A^{\otimes p}\otimes\Q(-p)$ by $(-1)^p(-1)^p=1$. This proves the complete assertion.

For a concrete counterexample take $X=\mathbf P^1_\R\times_\R\Spec\C$, a smooth projective real scheme. On $H^2(X_\C,\Q)$ with the two complex orientations, the algebraic deck involution has matrix
\[
 d=\begin{pmatrix}0&1\\1&0\end{pmatrix},\qquad
 f_\infty=-d.
\]
The algebraic projector $(1-d)/2$ has image $L$ and $f_\infty=1$ on that image. Its weight is $2$, whereas every tensor power of $\Q(2)$ has weight in $4\Z$. A Hodge morphism preserves weight, so the image has zero morphisms to and from every such power; it cannot be their sum. All cycles needed for this Artin--Tate example are algebraic, so the Hodge premise is satisfied here. The projector and the sign calculation are verified exactly in \texttt{code/verify\_exact.py}. This is the example recorded in RH10; we have checked its actual algebraic projector.
\end{proof}

\begin{result}{C.3.2}{Reality of the comparison point and of its torsor}{Proved}
For a motive over $\R$ in the preceding categories, use the right action on $P=\Isom^\otimes(T_{\B}\otimes\R,T_{\dR})$. Then
\[
 \bar\alpha=\alpha f_\infty,
 \qquad
 \alpha\in P(\R)\ \Longleftrightarrow\ f_\infty=1.
\]
The class of $P$ in $H^1(\R,G)$ is represented by $f_\infty$. If $\Q(-1)$ belongs to the generated category, $\R(\alpha)=\C$. Existence of some real point of $P$ does not imply that the distinguished $\alpha$ is real.
\end{result}
\begin{proof}
Transport the de Rham conjugation to Betti coordinates. It is $f_\infty c$, where $c$ conjugates coefficients, and $f_\infty$ is rational; hence it squares to one and its fixed space is $V^+_\R\oplus iV^-_\R$. This proves the comparison equation and the descent cocycle. Freeness of the torsor action gives the stated equivalence for $\alpha$. On $\Q(-1)$, conjugation reverses the orientation of $\mathbf P^1(\C)$, so $f_\infty=-1$. A nonreal coordinate of $\alpha$ together with $\R$ generates $\C$. On the Tate category, however, $i^{-1}\bar i=-1$ trivializes the torsor while the distinguished period coordinate is a nonzero real multiple of $2\pi i$ or its inverse, according to comparison direction. This verifies the exact distinction through a specified coboundary.
\end{proof}

\begin{result}{C.3.3}{Central real Frobenius}{Proved in the stated categories}
For a tensor generator $M$ in C.1.1, or under the hypothesis of C.1.2, the following conditions are equivalent: $f_\infty$ is central in $G_M$; every Hodge summand of $M$ has type $(p,p)$; $M$ belongs to the real Artin--Tate category generated by $A$ and $T=\Q(-1)$; $G_M$ is commutative.
\end{result}
\begin{proof}
Centrality and $f h(z)f^{-1}=h(\bar z)$ give $h(z)=h(\bar z)$. On type $(p,q)$ this says $z^{-p}\bar z^{-q}=\bar z^{-p}z^{-q}$ for all $z$, hence $p=q$. Full faithfulness and descent as in C.3.1 then express $M$ as a sum of $T^p$ and $A T^p$. Their tensor group is a quotient of $\mu_2\times\Gm$, so it is commutative. Commutativity implies centrality.

This also adjudicates the concrete candidate readings of the uncertain phrase ``$\Q$ tordu'' on 18:36. Reading it as $A$, the category $\langle A,\Q(2)\rangle$ misses $T$ by its weight. Reading it as $AT$, $\langle AT,\Q(2)\rangle=\langle AT\rangle$ has trivial $f$ and also misses $T$. Reading it as $T$ gives the Tate category and misses $A$. The precise correct generating category is $\langle A,T\rangle$. A phrase intended to denote the entire collection of Artin twists of Tate lines describes exactly that category; the supplied text does not specify that collective reading.
\end{proof}

\subsection{Transport, scalar fields and polarization}

\begin{result}{C.4.1}{Fixed-topology transport and the scalar fields on 18:41}{Proved}
For a topological field $K$ isomorphic to $\C$, the splitting of algebraic de Rham cohomology of a smooth projective $K$-variety obtained by Hodge theory is independent of the choice of topological isomorphism $K\simeq\C$, with coefficient conjugation included. For a field $k$ admitting a complex embedding and an archimedean place $v$, put $R_v=\{a\in k:v(a)\in\R\}$. The intersections defined on 18:41 are $k$-vector spaces for the Hodge summands and $R_v$-vector spaces for their Betti-real parts. Moreover
\[
 \bigcap_vR_v=\{a\in k:a\text{ is algebraic and totally real over }\Q\}.
\]
The intersection over all places equals the largest totally real subfield when $k$ is a number field.
\end{result}
\begin{proof}
The two continuous isomorphisms $K\to\C$ differ by complex conjugation. The comparison square uses the conjugate variety and the conjugate-linear maps both on de Rham coefficients and on singular-cohomology coefficients. In that square pullback by the antiholomorphic identification exchanges $(p,q)$ and $(q,p)$, while coefficient conjugation exchanges them again. Their composite preserves the labelled summands and the de Rham filtration. This is the required map behind the ``non comm.'' annotation on 18:39.

The summand in the completed cohomology is a vector space over $k_v$, so its intersection with $H_{\dR}(X/k)$ is closed under $k$. Intersecting also with the Betti-real form preserves multiplication by exactly the real scalars. A nonzero Betti-real subspace has scalar stabilizer $R_v$: a nonreal scalar times a nonzero real vector cannot be real, as comparison identifies this assertion with ordinary coordinate conjugation. Zero intersections impose no maximal-scalar claim.

Every totally real algebraic element lies in all $R_v$. An algebraic element with a nonreal conjugate is excluded by an embedding of its number field, extended to $k$. A transcendental element can be included in a transcendence basis of $k/\Qb\cap k$ and sent to a nonreal transcendental number, choosing the other basis elements algebraically independent and extending the resulting embedding through the algebraic extension. The cardinality needed for this choice is available because $k$ already embeds in $\C$. Thus no transcendental element belongs to the intersection. These are statements about the actual intersections on the page; they do not assert that their scalar extensions exhaust the completed Hodge summands.
\end{proof}

\begin{result}{C.5.1}{The positive Weil form}{Proved}
Let $M$ have pure weight $n$ and let $q:M\otimes M\to\Q(-n)$ be a polarization morphism over $\R$. Use the convention of \cite[(5.4)--(5.5)]{MilneShimura}, so $C=h(i)$ and
\[
 B(x,y)=(2\pi i)^n q(x,Cy)
\]
is symmetric positive definite on $H_\B(M)\otimes\R$. Then $fCf^{-1}=C^{-1}$ and $B(fx,fy)=B(x,y)$.
\end{result}
\begin{proof}
The normalizer identity with $i\mapsto-i$ gives the first equation. Since $q$ is a morphism over $\R$ and $f$ acts by $(-1)^n$ on its Tate target, $q(fx,fy)=(-1)^nq(x,y)$. Also $C^2=(-1)^n$. Consequently
\[
 \begin{aligned}
 B(fx,fy)
 &=(2\pi i)^nq(fx,fC^{-1}y)\\
 &=(-1)^n(2\pi i)^nq(x,C^{-1}y)\\
 &=(-1)^{2n}(2\pi i)^nq(x,Cy)=B(x,y).
 \end{aligned}
\]
Positivity is exactly the polarization convention just specified, rather than an additional assertion about every bilinear form. The weight-one sign identities have an exact matrix certificate in \texttt{verify\_exact.py}.
\end{proof}

\begin{result}{C.4.2}{The Hodge splitting over the real coefficient field}{Proved}
For a pure motive over $\R$ in the stated realization category, the de Rham Hodge splitting $\beta$ is $\R$-rational, and $\R(\gamma)=\R(\alpha)$. This is the coefficient-field assertion on 18:38.
\end{result}
\begin{proof}
Let $\kappa_{\B}$ be coefficient conjugation on $V_{\B,\C}$. Each of $\kappa_{\B}$ and $f_\infty$ exchanges the Hodge labels $(p,q)$ and $(q,p)$, so their commuting product preserves every labelled summand. Under comparison the real de Rham conjugation is $\kappa_{\dR}=\alpha f_\infty\kappa_{\B}\alpha^{-1}$. It therefore preserves each transported summand, and the isomorphism from a summand to its Hodge-graded quotient is defined over $\R$. Their direct sum is $\beta$, proving its $\R$-rationality. The identity $\gamma=(\beta\alpha)^{-1}$ and its inverse expression for $\alpha$ show equality of the two coordinate fields over $\R$. For an algebraic coefficient subfield $k\subset\R$, this calculation only places $\beta$ in $\R$; the separate question $\beta\in\overline{k}$ retains the gate in A.4c.
\end{proof}
